% ADDED figure for §CA.7 (2026-10-29); generated by qft-C/tools/ca7/figs/make_figs.sh
\tikzset{
  fdl/.style={line width=0.8pt, black},
  fdx/.style={circle, fill=black, inner sep=0pt, minimum size=3.6pt},
  fdv/.style={circle, fill=figB, draw=figB, inner sep=0pt, minimum size=5pt},
  leg/.style={circle, fill=figO, inner sep=0pt, minimum size=3.4pt},
  legn/.style={font=\tiny, text=figO!80!black, inner sep=1pt},
  vc/.style={draw=figB, densely dashed, line width=0.5pt},
  lab/.style={font=\scriptsize, inner sep=2pt},
  info/.style={font=\scriptsize, align=left, anchor=west},
  arr/.style={-{Stealth}, gray, line width=0.6pt},
  arrl/.style={font=\scriptsize, align=center, text=gray},
}
% \vtx{name}{center}{label}{angle/num,...}: dashed vertex circle, orange legs, leg numbers inside, name below
\newcommand{\vtx}[4]{%
  \draw[vc] (#2) circle (0.5); \node[font=\scriptsize, text=figB] at ($(#2)+(0,-0.72)$) {$#3$};
  \foreach \ang/\num in {#4}{%
    \coordinate (#1-\num) at ($(#2)+(\ang:0.5)$);
    \node[leg] at (#1-\num) {};
    \node[legn] at ($(#2)+(\ang:0.3)$) {\num};}}
% \chainlab{prefix}{left name}{right name}{left angle/label,...}{right angle/label,...}{xshift}
% chain of two tadpoles with leg slots a=180, b=125, c=55, d=0; blocks x1-a(L), b(L)-c(L), d(L)-a(R), b(R)-c(R), d(R)-x2
\newcommand{\chainlab}[6]{%
  \begin{scope}[xshift=#6]
    \coordinate (#1L) at (1.4,0); \coordinate (#1R) at (3.3,0);
    \draw[vc] (#1L) circle (0.45); \draw[vc] (#1R) circle (0.45);
    \node[font=\scriptsize, text=figB] at ($(#1L)+(0,-0.7)$) {$#2$};
    \node[font=\scriptsize, text=figB] at ($(#1R)+(0,-0.7)$) {$#3$};
    \foreach \s/\ang in {a/180,b/125,c/55,d/0}{%
      \coordinate (#1l\s) at ($(#1L)+(\ang:0.45)$); \coordinate (#1r\s) at ($(#1R)+(\ang:0.45)$);}
    \node[fdx] (#1x) at (0,0){}; \node[fdx] (#1y) at (4.7,0){};
    \node[lab,below] at (#1x){$x_1$}; \node[lab,below] at (#1y){$x_2$};
    \draw[fdl] (#1x)--(#1la); \draw[fdl] (#1ld)--(#1ra); \draw[fdl] (#1rd)--(#1y);
    \draw[fdl] (#1lb) .. controls ($(#1L)+(-0.55,1.45)$) and ($(#1L)+(0.55,1.45)$) .. (#1lc);
    \draw[fdl] (#1rb) .. controls ($(#1R)+(-0.55,1.45)$) and ($(#1R)+(0.55,1.45)$) .. (#1rc);
    \foreach \s in {a,b,c,d}{\node[leg] at (#1l\s){}; \node[leg] at (#1r\s){};}
    \foreach \ang/\t in {#4}{\node[legn] at ($(#1L)+(\ang:0.27)$){\t};}
    \foreach \ang/\t in {#5}{\node[legn] at ($(#1R)+(\ang:0.27)$){\t};}
  \end{scope}}
\begin{tikzpicture}[font=\small]
  \fill[gray!12, rounded corners=6pt] (3.15,-1.1) rectangle (5.25,1.0);
  \node[font=\scriptsize, text=gray!70!black] at (4.2,1.2) {not attached: $A(P) = \{2\}$};
  \coordinate (c1) at (1.5,0); \coordinate (c2) at (4.2,0);
  \vtx{a}{c1}{z_1}{180/1,125/2,55/3,0/4}
  \vtx{b}{c2}{z_2}{145/1,215/2,35/3,-35/4}
  \node[fdx] (x1) at (0,0){}; \node[fdx] (x2) at (3.0,0){};
  \node[lab,below] at (x1){$x_1$}; \node[lab,below] at (x2){$x_2$};
  \draw[fdl] (x1)--(a-1); \draw[fdl] (a-4)--(x2);
  \draw[fdl] (a-2) .. controls (0.95,1.5) and (2.05,1.5) .. (a-3);
  \draw[fdl] (b-1) .. controls (3.25,0.85) and (3.25,-0.85) .. (b-2);
  \draw[fdl] (b-3) .. controls (5.15,0.85) and (5.15,-0.85) .. (b-4);
  \node[font=\scriptsize] at (2.6,-1.45) {$P \in \mathcal P(X_{E,2})$, $E = \{x_1, x_2\}$};
  \draw[arr] (5.7,0.1)--(7.0,0.1) node[midway,above,arrl] {split};
  % P_att
  \begin{scope}[xshift=7.4cm]
    \coordinate (d1) at (1.5,0);
    \vtx{p}{d1}{z_1}{180/1,125/2,55/3,0/4}
    \node[fdx] (y1) at (0,0){}; \node[fdx] (y2) at (3.0,0){};
    \node[lab,below] at (y1){$x_1$}; \node[lab,below] at (y2){$x_2$};
    \draw[fdl] (y1)--(p-1); \draw[fdl] (p-4)--(y2);
    \draw[fdl] (p-2) .. controls (0.95,1.5) and (2.05,1.5) .. (p-3);
    \node[font=\scriptsize, align=center] at (1.5,-1.45) {$P_{\rm att}$: an attached pairing of\\ $E \sqcup (\{1\}\times[4])$};
  \end{scope}
  \node at (11.0,0) {$\times$};
  \begin{scope}[xshift=12.0cm]
    \coordinate (d2) at (1.0,0);
    \vtx{q}{d2}{z_2}{145/1,215/2,35/3,-35/4}
    \draw[fdl] (q-1) .. controls (0.05,0.85) and (0.05,-0.85) .. (q-2);
    \draw[fdl] (q-3) .. controls (1.95,0.85) and (1.95,-0.85) .. (q-4);
    \node[font=\scriptsize, align=center] at (1.0,-1.45) {$P_{\rm vac}$: a pairing of\\ $\{2\}\times[4]$};
  \end{scope}
  \node[font=\scriptsize, align=center] at (7.5,-2.5) {No block joins the two parts, so $P \mapsto (A(P), P_{\rm att}, P_{\rm vac})$ loses nothing and can be inverted: $P = P_{\rm att} \sqcup P_{\rm vac}$.};
\end{tikzpicture}
